\documentclass[../main.tex]{subfiles}
\graphicspath{{\subfix{../images/}}}
\begin{document}
	I circuiti in regime sinusoidale sono circuiti in corrente alternata le cui correnti e tensioni assumono un andamento sinusoidale, ovvero:
	\[e(t) = E_{MAX} \cdot \cos (\omega t + \phi)\]
	Dove:
	\begin{itemize}
		\item $E_{MAX}$ è l'ampiezza massima dell'onda
		\item $\omega$ è la pulsazione misurata in rad/s
		\item $\phi$ è la fase
	\end{itemize}
	Queste 3 grandezze che descrivono l’onda rimangono costanti durante tutto il regime sinusoidale.\\
	In particolare, si definiscono:
	\begin{itemize}
		\item La frequenza dell’onda come $f = \frac{\omega}{2 \pi} \;\; [Hz]$
		\item Il suo periodo come $T = \frac{1}{f} = \frac{2 \pi}{\omega} \;\; [s]$
		\item La sua pulsazione come $2 \pi f \;\; [rad/s]$
	\end{itemize}
	I circuiti in regimi sinusoidali sono anche detti circuiti isofrequenziali, proprio perché tutte le correnti e le tensioni coinvolte presentano la stessa pulsazione $\omega$.
	\subsection{Bipoli Elettrici}
	\subsubsection{Generatore Ideale di Tensione Sinusoidale}
	\subsubsection{Generatore Reale di Tensione Sinusoidale}
	\subsubsection{Generatore Ideale di Corrente Sinusoidale}
	\subsubsection{Generatore Reale di Corrente Sinusoidale}
	\subsection{Grandezze nel Regime Sinusoidale}
	Impedenza
	\subsection{Trasformata di Steinitz}
	\subsection{Metodo Risolutivo per Regime Sinusoidale}
	\subsection{Potenze in Regime Sinusoidale}
	\subsection{Valore Efficace di Tensioni e Correnti}
	\subsection{Teorema di Boucherot}
	\subsection{Rifasamento}
	\subsection{Teorema del Massimo Trasferimento di Potenza}
	\subsection{Sistemi Trifase}
	\subsubsection{Collegamento a Triangolo}
	\subsubsection{Collegamento a Stella}
	\subsubsection{Misura di Ptenza Trifase}
	\subsubsection{Inserzione Aaron}
\end{document}